% Chapter 5 — Evaluation and Discussion
% Plan mapping: PLANO §6.1 (row 5). Methodology of measurement is written now;
% ALL result slots are empty placeholders marked \todo{pending data-v1} and are
% filled only after the data freeze, from results/raw + manifests + the single
% analysis script, with entries in docs/claim_evidence_matrix.md.
% NO quantitative value may be added to this chapter by hand.
% Indicative target: 5,000–6,000 words.

\chapter{Evaluation and Discussion}
\label{ch:evaluation}

\noindent This chapter reports the experimental evaluation of the gateway
against the protocol frozen in Section~\ref{sec:method-protocol}. Every
quantitative statement in this chapter resolves to a raw-data directory with
its manifest and checksums, and is regenerated by the analysis pipeline; the
claim--evidence matrix in the repository records this mapping per claim.

\section{Measurement Methodology}
\label{sec:eval-method}

\noindent Latency is measured entirely within the controller process: a
monotonic-clock reading at the \ac{MQTT} receive callback and a second reading
after the Ditto acknowledgement, with the difference converted to
milliseconds. Using a single process and a monotonic clock avoids clock
synchronisation between machines for the primary metric. End-to-end
accounting (sent, received, duplicate, rejected, confirmed, lost) joins the
simulator's send log with the controller's event log on message identifiers,
which are deterministic per run. Resource consumption is sampled every second
per container. The statistical unit is the run; per-condition aggregates
report mean, standard deviation and 95\% confidence intervals, with medians
and percentiles for skewed distributions, as fixed in
Section~\ref{sec:method-stats}.

\section{Experimental Setup}
\label{sec:eval-setup}

\noindent \todo{pending data-v1 --- record from the campaign manifests: cloud
provider and region of the ARM64 \ac{VM}, exposed \ac{CPU} model and vCPU
count, memory and disk, kernel and operating-system versions, the shared-vCPU
limitation, container image digests, repository commit, and the exact
simulator host placement. No value may be written here that is not present in
a manifest under \texttt{results/raw/}.}

\section{Platform Reproducibility and Readiness}
\label{sec:eval-platform}

\noindent This section answers the functional part of RQ1.

\subsection{Image Build and Boot Consistency}
\todo{pending data-v1 --- clean-checkout build outcome, five QEMU boot
results, container-runtime smoke evidence.}

\subsection{Stack Time-to-Readiness}
\todo{pending data-v1 --- ten cold starts: readiness times with run-level
statistics.}

\subsection{Twin Creation}
\todo{pending data-v1 --- ten independent twin creations: outcomes.}

\section{Correctness and Reliability}
\label{sec:eval-correctness}

\noindent This section answers RQ2.

\subsection{Nominal Scenario}
\todo{pending data-v1 --- ten nominal runs: accepted/rejected/duplicate/lost
counts and delivery rate per run; final twin-state verification for the three
devices.}

\subsection{Fault Injection: Invalid Payloads}
\todo{pending data-v1 --- invalid-payload scenario: correct-rejection
accounting, no contamination of twins.}

\subsection{Fault Injection: Dropout and Reconnect}
\todo{pending data-v1 --- dropout-reconnect scenario: recovery behaviour,
duplicate handling after reconnect, event accounting.}

\section{Performance}
\label{sec:eval-performance}

\noindent This section answers RQ3.

\subsection{Latency}
\todo{pending data-v1 --- nominal latency p50/p95/p99 per run with run-level
aggregates; distribution figures generated by the analysis script.}

\subsection{Throughput and Saturation}
\todo{pending data-v1 --- load-sweep results at 10/50/100/250 msg/s; identify
the operational saturation point per the frozen definition (loss above 1\%,
persistent queue growth, p95 above one second, or sustained \ac{CPU} above
90\% for 60 seconds).}

\subsection{Resource Consumption}
\todo{pending data-v1 --- per-container \ac{CPU} and \ac{RAM} profiles across
load levels.}

\subsection{Stability: 24-Hour Soak}
\todo{pending data-v1 --- descriptive soak analysis: continuity of service,
resource trends, absence or presence of unrecovered failures. No confidence
interval is computed for the single soak run.}

\section{Answers to the Research Questions}
\label{sec:eval-answers}

\subsection{RQ1}
\todo{pending data-v1 --- answer RQ1 strictly from the platform evidence
above.}

\subsection{RQ2}
\todo{pending data-v1 --- answer RQ2 strictly from the correctness and
reliability evidence above.}

\subsection{RQ3}
\todo{pending data-v1 --- answer RQ3 strictly from the performance evidence
above, stating the observed trade-offs and their operational envelope.}

\section{Comparison with Prior Work}
\label{sec:eval-comparison}

\noindent \todo{pending data-v1 --- qualitative comparison of the evaluated
scope against the published C2DTA evaluation \cite{pinto2026c2dta} (single
smartwatch profile, x86 virtual machine) and against edge digital-twin
evaluations such as \cite{picone2023flexible}; strictly no numeric
cross-paper comparisons unless measurement conditions are demonstrably
comparable.}

\section{Threats to Validity}
\label{sec:eval-validity}

\noindent The following threats are identified at protocol time; the
discussion is revisited once results exist.

\begin{itemize}
  \item \textbf{Internal validity.} Instrumentation shares the process with
        the measured code; the event-log write path adds work per message.
        Mitigations: monotonic clocks, single-process latency measurement,
        warm-up exclusion, and identical instrumentation across all
        conditions. The cloud \ac{VM}'s shared vCPUs may introduce
        neighbour-induced variance; mitigations: repeated runs, randomised
        load-sweep order, and reporting of dispersion, with the limitation
        stated explicitly.
  \item \textbf{External validity.} Telemetry is synthetic and the deployment
        is a single gateway with three devices on one ARM64 cloud \ac{VM};
        results do not automatically transfer to physical single-board
        hardware, other device populations or other stacks. Emulated-platform
        evidence is functional only.
  \item \textbf{Construct validity.} The primary latency metric covers the
        controller's processing path (receive to twin acknowledgement), not
        the device-to-gateway network leg; delivery-rate definitions count
        deliberately invalid payloads as rejections, not losses. These
        definitions are fixed in Section~\ref{sec:method-stats} and applied
        uniformly.
  \item \textbf{Conclusion validity.} The statistical unit is the run;
        confidence intervals use the per-run distribution, and the single
        soak run is analysed descriptively without inferential statistics.
\end{itemize}
